\documentclass[../DoAn.tex]{subfiles}
\begin{document}

\begin{center}
    \Large{\textbf{TÓM TẮT NỘI DUNG ĐỒ ÁN}}\\
\end{center}
\vspace{1cm}

Theo dõi vị trí trẻ em theo thời gian thực là một bài toán thực tế xuất phát từ nhu cầu an toàn của các gia đình hiện đại. Các giải pháp hiện có trên thị trường đều có hạn chế đáng kể: ứng dụng Life360 phụ thuộc vào điện thoại thông minh của trẻ, thiết bị AirTag chỉ hoạt động trong phạm vi ngắn và không hỗ trợ theo dõi liên tục, còn các GPS tracker thương mại thường đi kèm phí đăng ký dịch vụ cao và thiếu khả năng tùy chỉnh. Từ phân tích đó, đồ án đặt mục tiêu xây dựng một hệ thống hoạt động độc lập, không phụ thuộc vào thiết bị di động của trẻ, với chi phí vận hành thấp và khả năng tùy biến linh hoạt theo nhu cầu thực tế.

Hệ thống được thiết kế theo kiến trúc ba tầng gồm thiết bị phần cứng IoT, backend xử lý và ứng dụng di động dành cho phụ huynh. Phần cứng sử dụng vi điều khiển ESP32 kết hợp với module SIM7600 -- một linh kiện tích hợp đồng thời bộ thu GPS và modem 4G LTE trong một vỏ duy nhất, cho phép thiết bị định vị và truyền dữ liệu mà không cần thêm module ngoài. Dữ liệu vị trí được truyền định kỳ qua giao thức MQTT lên backend xây dựng bằng FastAPI, lưu trữ trong MongoDB và phân phối đến ứng dụng Flutter qua WebSocket.

Đóng góp kỹ thuật chính của đề tài tập trung vào bốn điểm. Thứ nhất, tích hợp GPS và kết nối di động trên cùng một module kết hợp với firmware máy trạng thái có khả năng tự phục hồi khi mất tín hiệu. Thứ hai, thuật toán geofencing hai chế độ hỗ trợ cả vùng tròn theo công thức Haversine và hành lang đường đi theo phép chiếu điểm-đoạn thẳng. Thứ ba, kiến trúc truyền thông lai MQTT và WebSocket sử dụng hàng đợi bất đồng bộ để chuyển tiếp bản tin giữa hai giao thức mà không gây khóa luồng. Thứ tư, cơ chế cảnh báo đa kênh song song qua WebSocket, Telegram Bot và Email với thời gian chờ chống gửi trùng.

Kết quả đo thực tế cho thấy độ trễ end-to-end từ thiết bị đến ứng dụng dưới 2 giây, độ chính xác GPS khoảng 3,2 m ngoài trời và chu kỳ cập nhật 5 giây. Hệ thống đã vận hành ổn định trong điều kiện thực tế và đáp ứng đầy đủ các yêu cầu đặt ra.

\begin{flushright}
Sinh viên thực hiện\\
\vspace{0.5cm}
Nguyễn Đức Dương
\end{flushright}

\end{document}
